\section{Introduction}
\label{sec:introduction}

Learning-based controllers have enabled major advances in robotics.
%
Much of this progress relies on online reinforcement learning (RL),
which often requires millions of environment interactions to learn a
successful policy.
%
Massively parallel simulation generates this experience efficiently by
running thousands of robots at the same
time~\cite{makoviychuk2021isaacgym,mittal2025isaaclab}.
%
This approach is particularly valuable for aerial and marine robots,
where trial-and-error learning on hardware risks crashes, requires
costly deployments, and encounters conditions that are difficult to
repeat~\cite{grimaldi2025stonefish}.
%
Simulation provides a controlled setting for training policies before
deployment.

Training these robots requires simulation that is both fast and
accurate enough to capture the dynamics relevant to the task.
%
Detailed actuator and fluid models can capture these dynamics, but
their computational cost can make large-scale parallel training
impractical.
%
Simpler models run faster, but may omit physical effects that matter
when learned policies are deployed on real robots.
%
The appropriate balance depends on the vehicle and task.
%
Supporting different levels of physical detail could therefore make
policy learning practical across a wider range of aerial and marine
applications.

Many existing aerial and marine simulators rely on CPU-based physics,
which can limit how efficiently they scale to large numbers of
training environments.
%
Stonefish provides detailed marine dynamics using CPU-based
simulation~\cite{grimaldi2025stonefish}, while Flightmare runs quadrotor
dynamics in parallel across CPU
threads~\cite{song2021flightmare}.
%
GPU-based simulators such as Aerial Gym support massively parallel
training, but focus primarily on multirotor
platforms~\cite{kulkarni2025aerial} while relying on a simplified model.
%
Extending these benefits to aircraft, surface vessels, and underwater
vehicles requires a shared framework that combines GPU-parallel
simulation with reusable actuator and fluid models.

Isaac Lab and Newton provide a foundation for such a
framework~\cite{mittal2025isaaclab,newton2025}.
%
They support joint actuation as well as external forces and torques
applied to robot bodies.
%
However, applying forces is only one part of integrating a physical
model into the simulation.
%
The model also needs commands, body motion, and fluid conditions as
inputs.
%
Its internal state must advance with the simulation and reset when
the corresponding environment resets.
%
When each task handles this logic separately, moving a model to another
vehicle requires rewriting the code that connects it to the simulation.

We present \textbf{aFluid}, a modular framework for learning to control
robots in fluids.
%
Built on Newton, aFluid provides a library of analytical actuator and
fluid models, a shared GPU runtime, and an Isaac Lab integration.
%
The runtime executes the models across parallel environments and
handles their state updates and resets.
%
Models can be selected, replaced, or extended for each vehicle while
retaining the same tools for training, randomization, and evaluation.

The key idea is to connect physical models through their common
output: forces and torques on a body.
%
Each propeller, lifting-surface, or buoyancy model contributes a
wrench, expressed in a common coordinate frame and referenced to the
body's center of mass.
%
These contributions are added together and passed to Newton to
advance the robot's motion.
%
Fluid conditions are defined separately from the force models,
allowing the same model implementation to be used across vehicles
and media with appropriate parameters.
%
For example, a lifting-surface model can describe a sail in air or a
rudder in water using the corresponding fluid properties and force
coefficients.
%
The physical models can therefore change without changing the
surrounding simulation and training workflow.

In preliminary simulation experiments, we train policies for quadrotor
control, flips with recovery, assisted fixed-wing navigation, and
sailboat control in up to 4,096 parallel environments.
%
In evaluations with new random seeds, the policies complete 254 of
256 flip-and-recovery maneuvers and all 256 fixed-wing waypoint
courses.
%
The sailboat policy completes 232 of 256 courses under randomized
steady wind.
%
All tasks use the same simulation and training interfaces while
combining different actuator and fluid models.